\chapter{REPRODUTIBILIDADE COMPUTACIONAL}
\label{ap:reprodutibilidade}

\section{Repositório e versão auditada}

O código-fonte está disponível em:

\begin{center}
\url{https://github.com/paulorenatoaz/coinfosim}
\end{center}

Este relatório documenta o \CoInfoSim{} \CoInfoSimVersion{}. Os quatro indicadores de preservação do perfil de cooperação preditiva ($\rankrho$, $\AW$, $\AR$, $\DR$, $\SR$) foram computados diretamente pelo código atual (\texttt{coinfosim.results.predictive\_profile}) para os cenários de escala completa (Capítulo~5), a partir dos resultados brutos por replicação persistidos, sem nova simulação de Monte Carlo. A citação do software é fornecida em \texttt{CITATION.cff} e na entrada bibliográfica \texttt{azevedo2026coinfosim} \cite{azevedo2026coinfosim}.

\section{Estrutura dos artefatos persistidos}

Cada execução de simulação persiste configuração, metadados, estado do acumulador, resumo de parada e dados necessários para regenerar os relatórios. Os registros de cada execução incluem:

\begin{itemize}
  \item tamanhos amostrais e subconjuntos avaliados;
  \item classificadores e parâmetros resolvidos;
  \item tabelas de perda média estimada e erro-padrão;
  \item rankings por $n$;
  \item matriz de vencedores $\Wmat$ e eventos de inversão de vencedor;
  \item indicadores progressivos de preservação do perfil ($\rankrho$, $\AW$, $\AR$, $\DR$, $\SR$), esquema \texttt{schema\_version=3};
  \item tamanho e origem do conjunto de teste real fixo;
  \item tempo de execução e política de paralelização;
  \item número de replicações e razão de parada para cada $n$.
\end{itemize}

Os arquivos \texttt{scenario.json} e \texttt{simulation.json} contêm registros legíveis; o resultado completo permanece em arquivos compactados adjacentes para a regeneração exata dos relatórios. Cada cenário regenerado emite ainda, ao lado do relatório HTML, o manifesto semântico e os artefatos de proveniência formal (\texttt{semantic\_manifest.json}, \texttt{provenance.provjson}, \texttt{provenance.provn}, \texttt{provenance.ttl}); este apêndice descreve apenas onde e como esses artefatos são produzidos --- sua interpretação formal, o mapeamento de tipos PROV-O e os limites da camada semântica são detalhados no Apêndice~E.

\section{Configuração experimental comum}

Nos cenários completos, a grade utilizada foi
\[
\{2,4,8,16,32,64,128,256,512\}
\]
amostras por classe. A regra de parada empregou mínimo de 100 replicações, máximo de 2000, lotes de 20 e alvo de meia-largura de intervalo de confiança de 95\% igual a 0,01. A parada ocorre quando a maior meia-largura entre células $(S_I,f)$ atinge a tolerância ou quando o orçamento máximo é consumido.

A padronização utiliza momentos estimados apenas nos dados reais de treinamento e $\mathrm{ddof}=0$. O mesmo conjunto de teste real fixo é reutilizado nas três condições experimentais. As amostras são balanceadas por classe e incrementalmente aninhadas dentro de cada replicação.

\section{Política de sementes}

Cada execução registra a semente-base e os identificadores derivados. Para os classificadores determinísticos ou historicamente fixados, os parâmetros de aleatoriedade permanecem congelados. No Random Forest, a política versionada \texttt{classifier\_seed\_v1} deriva a semente do estimador a partir da semente-base, do espaço de nomes do classificador e do ID da replicação. A mesma semente de replicação é usada entre condições experimentais, tamanhos amostrais e subconjuntos.

Essa política implementa uma comparação pareada. Para estudos de sensibilidade, recomenda-se executar uma política alternativa com sementes independentes entre condições e comparar as métricas estruturais.

\section{Calibração do Random Forest}

A calibração é realizada uma única vez com os dados reais de treinamento do SUPPORT2. O artefato versionado é
\begin{center}
\texttt{config/calibration/support2\_random\_forest.json}.
\end{center}

O comando documentado é:

\begin{verbatim}
.venv/bin/python scripts/calibrate_support2_random_forest.py \
  --raw-dir data/raw/support2 \
  --output config/calibration/support2_random_forest.json \
  --calibration-seed 0
\end{verbatim}

O artefato registra o hash do arquivo bruto, a impressão digital da partição de treinamento, a variável-alvo, a ordem dos canais, o particionamento, os parâmetros do estimador e a versão do scikit-learn. A execução normal não recalibra o modelo nem adota uma configuração alternativa silenciosa quando o artefato é inválido.

\section{Paralelização}

O paralelismo ocorre no nível externo das replicações Monte Carlo. Cada processo executa o estimador com paralelismo interno limitado a uma thread; no Random Forest, \texttt{n\_jobs=1} é obrigatório. Essa estratégia evita oversubscription e torna o uso de CPU previsível.

Um exemplo de execução de validação é:

\begin{verbatim}
.venv/bin/python scripts/run_support2_scenario.py \
  --mode smoke \
  --raw-dir data/raw/support2 \
  --output-dir output/validation/support2-random-forest-smoke \
  --rf-calibration-file \
  config/calibration/support2_random_forest.json \
  --execution-backend process \
  --n-jobs 5 \
  --worker-inner-threads 1 \
  --multiprocessing-start-method forkserver \
  --no-color
\end{verbatim}

O número de processos de trabalho deve ser ajustado à máquina, sem alterar as sementes ou a configuração científica.

\section{Regeneração dos relatórios}

Os relatórios podem ser regenerados a partir dos resultados persistidos sem nova execução Monte Carlo, por meio de:

\begin{verbatim}
coinfosim scenario regenerate <slug> --run-id <id> \
  --output-dir <diretório do registro de execução>
\end{verbatim}

Essa separação mantém o cálculo científico independente da apresentação. A publicação no GitHub Pages copia os diretórios de relatórios e dados para o ramo de publicação e reconstrói a página inicial que aponta para os cenários disponíveis.

\begin{limitationbox}[Terminologia do renderizador HTML]
O renderizador de cenário e a camada de análise (\texttt{coinfosim.results.predictive\_profile}) usam a mesma terminologia canônica --- $\Wmat$, $\Rmat$, $\rankrho$, $\AW$, $\AR$, $\SR$, $\DR$ --- em \texttt{schema\_version=3}, idêntica à deste relatório.
\end{limitationbox}

\section{Regeneração direta dos indicadores $\Wmat$/$\Rmat$ a partir de resultados persistidos}
\label{ap:comando-wr}

Para obter indicadores de preservação genuínos e correntes sem modificar código de produção nem executar Monte Carlo, esta revisão chamou diretamente as funções canônicas de \texttt{coinfosim.results.predictive\_profile} sobre os resultados brutos por replicação regenerados dos quatro cenários de escala completa (Occupancy 000002, Air Quality 000005, SUPPORT2 000007 e 000008), conforme a política \texttt{structural\_snapshot\_policy} de cada cenário (Apêndice~A): para 000002, 000005 e 000007 (\texttt{embedded}), os quatro indicadores foram lidos diretamente do payload \texttt{predictive\_cooperation\_profile} já persistido em \texttt{scenario.json}; para 000008 (\texttt{regenerate\_from\_result\_data}), os indicadores foram recalculados diretamente sobre os resultados brutos das simulações 000024--000026 em \texttt{output/reports/simulations/} (Apêndice~A):

\begin{verbatim}
.venv/bin/python3 -c "
from coinfosim.results.persistence import (
    load_simulation_result,
)
from coinfosim.results.predictive_profile import (
    scenario_predictive_profile_agreement,
)

base = 'output/reports/simulations'
arm_dirs = {
    'real_to_real': '000024_support2_real_data_full',
    'single_gaussian_to_real':
        '000025_support2_single_gaussian_to_real_full',
    'gmm_to_real': '000026_support2_gmm_to_real_full',
}
suffixes = {
    'real_to_real': '000024',
    'single_gaussian_to_real': '000025',
    'gmm_to_real': '000026',
}
paths = {
    arm: f'{base}/{d}/result_data_full_{suffixes[arm]}.json.gz'
    for arm, d in arm_dirs.items()
}
results = {
    arm: load_simulation_result(p) for arm, p in paths.items()
}
labels = {'real_to_real': 'Real -> Real',
          'single_gaussian_to_real': 'Single Gaussian -> Real',
          'gmm_to_real': 'GMM -> Real'}
profile = scenario_predictive_profile_agreement(
    results, 'real_to_real', labels
)
"
\end{verbatim}

O mesmo padrão, combinado com as funções do módulo
\path{coinfosim.reports.predictive_profile_visualization},

\begin{itemize}
  \item \texttt{winner\_matrix\_figure};
  \item \texttt{reversal\_matrix\_figure};
  \item \texttt{profile\_metric\_series\_figure},
\end{itemize}

\noindent produziu, respectivamente, as matrizes $\Wmat$/$\Rmat$ em \path{figures/wr_matrices/support2_full/} (Apêndice~D) e as figuras de evolução dos quatro indicadores por $n$ em \path{figures/predictive_profile_metrics/} (Capítulo~5). Este comando é reprodutível por qualquer pessoa com acesso ao repositório e ao diretório \path{output/reports/} (rastreado fora do controle de versão, mas presente no ambiente auditado); ele não requer, nem executa, nenhuma nova simulação. Nenhum resultado de execução \textit{smoke} foi utilizado para gerar figuras ou tabelas deste relatório; o modo \textit{smoke} do \CoInfoSim{} (Seção anterior) é reservado a testes e validação computacional do código, não a produção de evidência científica.

\section{Compilação deste relatório}

Este documento usa XeLaTeX e Biber. A compilação recomendada é:

\begin{verbatim}
latexmk -xelatex main.tex
\end{verbatim}

O preâmbulo segue o manual acadêmico da UFRJ para página A4, margens, espaçamento, paginação e estrutura, com elementos tipográficos inspirados no template da Sociedade Brasileira de Computação. O PDF final é renderizado e inspecionado visualmente após qualquer alteração de figuras, tabelas ou metadados, incluindo ao final desta revisão.

\section{Checklist mínimo de replicação}

Uma replicação independente deve registrar:

\begin{enumerate}
  \item commit do código e hash dos dados brutos;
  \item versão do Python e das dependências;
  \item comando completo de execução;
  \item sementes e política de derivação;
  \item particionamento e parâmetros de padronização;
  \item configuração dos geradores e classificadores;
  \item mecanismo de execução, processos de trabalho e linhas de execução internas;
  \item critérios de parada e replicações efetivas;
  \item IDs de simulação e cenário;
  \item hashes dos artefatos e relatórios publicados.
\end{enumerate}
